\section{API REST y endpoints}

Este capítulo documenta el contrato REST del nodo: cada endpoint con su método,
ruta, cuerpo, respuesta y códigos de estado. Toda la API se sirve sobre el mismo
router (\codigo{net/Routes.java}) y dos de los endpoints exigen un JWT,
siguiendo la estructura REST + JWT que pide la rúbrica.

\subsection{Mapa de rutas}

El enrutador vive en \codigo{net/Routes.java}. Es un \emph{router} por método y
ruta: cada patrón se parte en segmentos separados por \codigo{/} y un segmento
escrito como \codigo{\{nombre\}} es una variable de un solo segmento; el resto
debe coincidir por igualdad exacta de cadena (no es coincidencia por prefijo más
largo). El método \codigo{Routes.match()} resuelve un método y una ruta concreta
contra las rutas registradas con una semántica clara de tres resultados: si el
método y la forma de los segmentos coinciden devuelve el \codigo{Endpoint} con
estado \codigo{200}; si algún patrón coincide con la ruta pero ningún método
coincide devuelve \codigo{405} (\emph{method not allowed}); y si ningún patrón
coincide con la ruta devuelve \codigo{404}.

El cableado real de rutas a \emph{endpoints} ocurre en \codigo{app/Node.java},
mediante llamadas a \codigo{routes.register(...)}. La siguiente tabla resume cada
endpoint registrado, con su método, ruta, si exige autenticación con \emph{token}
y su propósito.

\begin{center}
\begin{tabular}{l l c l}
\textbf{Método} & \textbf{Ruta} & \textbf{Auth} & \textbf{Propósito} \\
\hline
POST & \codigo{/api/register} & No & Crea una credencial de usuario. \\
POST & \codigo{/api/login} & No & Verifica credencial y emite un JWT. \\
GET & \codigo{/api/accounts/\{id\}} & Sí & Devuelve el saldo de una cuenta. \\
POST & \codigo{/api/transactions/transfer} & Sí & Transfiere dinero entre cuentas. \\
GET & \codigo{/api/stats} & No & Métricas de \emph{este} nodo. \\
GET & \codigo{/panel} & No & Métricas agregadas del clúster. \\
GET & \codigo{/} & No & Sirve el tablero de una sola página. \\
\end{tabular}
\end{center}

Solo los dos endpoints de negocio (\codigo{/api/accounts/\{id\}} y
\codigo{/api/transactions/transfer}) exigen un \emph{token} válido; los de
autenticación y los de monitoreo son públicos. La columna \emph{Auth} se valida
con \codigo{Authenticator.authorize()} dentro de cada endpoint, no en el router.

\subsection{Autenticacion: register y login}

\subsubsection*{Registro de usuario}

\codigo{api/RegisterEndpoint.java} atiende \codigo{POST /api/register} y crea una
credencial nueva. Acepta un cuerpo JSON con la forma
\codigo{\{"username":"...","password":"..."\}}, lo parsea con \codigo{Gson} y
delega la persistencia en \codigo{Authenticator.signup()}, que cifra la
contraseña y la guarda en el \codigo{CredentialStore}. El endpoint primero
verifica el método (responde \codigo{405} si no es \codigo{POST}); si el JSON está
mal formado o falta \codigo{username}/\codigo{password} responde \codigo{400} con
\codigo{\{"error":"bad\_request"\}}; si el usuario aún no existía devuelve
\codigo{201} con \codigo{\{"status":"created"\}}; y si el nombre ya estaba tomado
devuelve \codigo{409} con \codigo{\{"error":"username\_taken"\}}. La distinción
entre creado (\codigo{201}) y duplicado (\codigo{409}) la decide directamente el
valor booleano que regresa \codigo{auth.signup(username, password)}.

\subsubsection*{Inicio de sesión}

\codigo{api/LoginEndpoint.java} atiende \codigo{POST /api/login}: verifica la
credencial y, en caso de éxito, emite un JWT. Recibe el mismo cuerpo
\codigo{\{"username":"...","password":"..."\}} y delega en
\codigo{Authenticator.login()}, que busca la credencial en el
\codigo{CredentialStore}, compara la contraseña con \codigo{Passwords.matches()} y
solo si coinciden mina un \emph{token} firmado con \codigo{Tokens.issue()}. Es
decir, el login valida credenciales reales: si el usuario no existe o la
contraseña no coincide, \codigo{login()} regresa \codigo{null} y el endpoint
responde \codigo{401} con \codigo{\{"error":"invalid\_credentials"\}}. Con
credenciales válidas responde \codigo{200} con \codigo{\{"token":"<jwt>"\}}; un
JSON mal formado produce \codigo{400} y un método distinto de \codigo{POST}
produce \codigo{405}. El cliente reutiliza ese \emph{token} en la cabecera
\codigo{Authorization: Bearer <jwt>} para las operaciones protegidas.

\begin{lstlisting}[language=none,caption={Contrato JSON de POST /api/login (peticion y respuesta)}]
POST /api/login
Content-Type: application/json

{"username":"juan","password":"secreto"}

200 OK
{"token":"<jwt>"}

401 Unauthorized
{"error":"invalid_credentials"}
\end{lstlisting}

\subsection{Operaciones de negocio: balance y transfer}

\subsubsection*{Consulta de saldo}

\codigo{api/BalanceEndpoint.java} atiende \codigo{GET /api/accounts/\{id\}} y
devuelve la instantánea de una cuenta. Exige un JWT válido: lo primero que hace,
tras descartar métodos distintos de \codigo{GET} (\codigo{405}), es llamar a
\codigo{auth.authorize(req)}; si regresa \codigo{null} responde \codigo{401}.
Luego resuelve la variable de ruta \codigo{id} con \codigo{req.pathParam("id")}:
si no es un entero (no encaja en el patrón \codigo{-?\textbackslash{}d+})
responde \codigo{400} con \codigo{\{"error":"bad\_id"\}}; si es numérico pero se
sale del rango de \codigo{int}, o si la cuenta no está registrada, responde
\codigo{404}. Con una cuenta válida construye el cuerpo con la forma obligatoria
del contrato, donde las llaves son exactamente \codigo{id}, \codigo{propietario}
y \codigo{balance}, por ejemplo
\codigo{\{"id":125,"propietario":"JUAN MOLINAR HERNANDEZ","balance":15750.25\}}.
Un detalle clave: \codigo{id} y \codigo{balance} deben serializarse como
\emph{números} JSON (sin comillas); como \codigo{Money.toDecimal()} devuelve un
\codigo{String}, el saldo se envuelve en un \codigo{BigDecimal} para que
\codigo{Gson} lo emita como número y no como cadena entre comillas.

\subsubsection*{Transferencia}

\codigo{api/TransferEndpoint.java} atiende \codigo{POST /api/transactions/transfer}
y mueve dinero entre dos cuentas. También exige JWT (\codigo{401} sin
\emph{token}) y rechaza métodos distintos de \codigo{POST} (\codigo{405}). El
cuerpo obligatorio usa la forma
\codigo{\{"sourceAccountId":"123","targetAccountId":"456","amount":200.00\}},
donde los identificadores son \emph{cadenas} y \codigo{amount} es decimal. El
monto decimal se convierte a centavos con \codigo{Money.toCents()} y la operación
se orquesta en \codigo{Bank.transfer()}, que valida la solicitud, aplica el
movimiento atómico, estampa un número de secuencia fresco y preserva la
invariante de conservación. Un éxito devuelve \codigo{200} con
\codigo{\{"status":"ok","seq":<seq>\}} llevando la secuencia asignada. Un rechazo
se manifiesta como una \codigo{TransferException}: su \codigo{code()} determina el
estado, donde \codigo{"no\_such\_account"} mapea a \codigo{404} y el resto de
códigos (\codigo{"self\_transfer"}, \codigo{"bad\_amount"}, \codigo{"low\_balance"})
mapean a \codigo{400}, y el cuerpo refleja \codigo{\{"error":<code>\}}. Un JSON
mal formado, un identificador no numérico o un monto inválido producen
\codigo{400} con \codigo{\{"error":"bad\_request"\}}.

\begin{lstlisting}[language=none,caption={Contrato JSON de POST /api/transactions/transfer}]
POST /api/transactions/transfer
Authorization: Bearer <jwt>
Content-Type: application/json

{"sourceAccountId":"123","targetAccountId":"456","amount":200.00}

200 OK
{"status":"ok","seq":42}

400 Bad Request
{"error":"low_balance"}

404 Not Found
{"error":"no_such_account"}
\end{lstlisting}

El mapeo de \codigo{TransferException} a estado HTTP se concentra en pocas
líneas, copiadas tal cual del endpoint:

\begin{lstlisting}[language=Java,caption={Mapeo de rechazos de transferencia a estado REST en TransferEndpoint.java}]
} catch (TransferException e) {
    int status = "no_such_account".equals(e.code()) ? 404 : 400;
    JsonObject err = new JsonObject();
    err.addProperty("error", e.code());
    return Reply.json(status, gson.toJson(err));
}
\end{lstlisting}

\subsection{Endpoints de monitoreo}

Tres endpoints alimentan el tablero. Aquí se resumen; su funcionamiento detallado
se trata en el capítulo 9.

\codigo{api/StatsEndpoint.java} atiende \codigo{GET /api/stats} y reporta las
métricas de \emph{este} nodo. Serializa la instantánea viva de
\codigo{NodeStats} llamando a \codigo{stats.toJson()}, que produce un objeto JSON
con las llaves \codigo{nodeId}, \codigo{role}, \codigo{status},
\codigo{accountCount}, \codigo{totalBalance}, \codigo{transferCount},
\codigo{lastTxId}, \codigo{cpuPercent}, \codigo{ramPercent}, \codigo{diskPercent}
y \codigo{gcsCount} (el número de transacciones almacenadas en Cloud Storage).
Responde \codigo{405} ante un método distinto de \codigo{GET}.

\codigo{api/PanelEndpoint.java} atiende \codigo{GET /panel} y agrega las métricas
de todo el clúster. Combina la instantánea local (\codigo{stats.toJson()}) con las
métricas que obtiene por HTTP del endpoint \codigo{/api/stats} de cada par
listado en \codigo{TES\_PEERS} (vía \codigo{NodeConfig}), y produce un único
documento JSON con la forma \codigo{\{"nodes":[...]\}}, una entrada por nodo, para
que el tablero dibuje el clúster en una sola página. Cada par se consulta de forma
concurrente con un \codigo{HttpClient} de \emph{timeout} acotado; un par
inalcanzable o lento no tumba la respuesta: se reporta como nodo caído
(\codigo{"status":"down"}) en lugar de fallar todo el panel. Además recuerda el
último \codigo{nodeId} que cada par reportó mientras estuvo arriba, de modo que la
tarjeta de un nodo caído conserve su identidad real durante la demostración de
tolerancia a fallos.

\codigo{api/DashboardEndpoint.java} atiende \codigo{GET /} y sirve el tablero de
una sola página: transmite el recurso \codigo{dashboard.html} desde el
\emph{classpath} como respuesta HTML. La propia página sondea \codigo{/panel} (y
ofrece un control de \emph{Refresh}) para mostrar el estado de cada nodo, su
número de cuentas, el saldo total, el número de transferencias, el id de la
última transacción, \%CPU, \%RAM, \%Disco y el conteo de transacciones en Cloud
Storage. Responde \codigo{405} ante un método distinto de \codigo{GET} y
\codigo{404} si el recurso HTML no existe.
